\title{DeepSeekMath: Pushing the Limits of Mathematical Reasoning in Open Language Models}

\begin{abstract}
The inherent complexity and structured requirements of mathematical reasoning present notable obstacles for language models. This study presents DeepSeekMath~7B, which extends the pre-training of DeepSeek-Coder-Base-v1.5 7B by incorporating 120B mathematics-focused tokens, drawing from Common Crawl along with datasets containing natural language and code. DeepSeekMath~7B attains a remarkable 51.7\% on the MATH benchmark—known for its competition-level problems—without the use of auxiliary toolkits or voting mechanisms, nearing the capabilities of leading models like Gemini-Ultra and GPT-4. Additionally, a self-consistency method leveraging 64 responses from DeepSeekMath~7B elevates performance to 60.9\% on MATH. The strong mathematical reasoning performance observed in DeepSeekMath~is attributable to two main strategies: firstly, the use of a carefully crafted data filtering pipeline to fully exploit large-scale, public web data; and secondly, the introduction of Group Relative Policy Optimization (GRPO), an adaptation of Proximal Policy Optimization (PPO), which not only advances mathematical reasoning but also improves the memory efficiency of PPO.
\end{abstract}

\section{Introduction}

Large language models (LLM) have transformed methods for addressing mathematical reasoning in artificial intelligence, driving forward progress on benchmarks for quantitative reasoning~\citep{MATH} and geometric reasoning~\citep{trinh2024solving}. These systems have also been highly beneficial in supporting human efforts to solve intricate mathematical problems~\citep{tao}. Despite these advancements, top-tier systems such as GPT-4~\citep{gpt4} and Gemini-Ultra~\citep{gemini} remain closed-source, leaving the available open-source alternatives with a substantial gap in capability and performance.

This paper presents DeepSeekMath, a language model specifically tailored for mathematical domains that markedly surpasses the mathematical proficiency of existing open-source models, achieving performance levels close to those of GPT-4 on established academic tasks. Central to this advancement is the construction of the DeepSeekMath Corpus, an extensive and meticulously curated pre-training dataset consisting of 120B tokens drawn from mathematical texts. The corpus is assembled by leveraging a fastText-based classifier \citep{joulin2016fasttext} on data sourced from the Common Crawl (CC). Initially, the classifier is trained using positive samples from OpenWebMath \citep{openwebmath}, contrasted against a varied sample of non-mathematical web content as negatives. The trained model is then used to identify additional candidate documents from the CC, which are subsequently filtered and validated by human annotators. These newly vetted instances are incorporated into the training set, enabling further refinement of the classifier. Experimental outcomes demonstrate that the resulting large-scale dataset possesses high fidelity: DeepSeekMath-Base 7B records a 64.2\% score on GSM8K \citep{gsm8k} and 36.2\% on the advanced MATH benchmark \citep{MATH}, thereby surpassing Minerva 540B \citep{minerva}. Owing to the multilingual design of the DeepSeekMath Corpus, we also observe notable gains in Chinese mathematical assessment benchmarks \citep{wei2023cmath,agieval}. Our methodology for compiling mathematical data aims to provide a foundation for the broader research community, with substantial opportunities remaining for further enhancement.

DeepSeekMath-Base is constructed upon DeepSeek-Coder-Base-v1.5 7B \citep{deepseek-coder}, as preliminary experiments reveal that code-centric pretraining provides a superior foundation relative to standard general-purpose language models. Additionally, we observe that incorporating mathematical pretraining not only bolsters mathematical competence but also yields notable gains on the MMLU \citep{mmlu} and BBH \citep{bbh} benchmarks, suggesting an enhancement in broader reasoning proficiency.

Post pre-training, we subject DeepSeekMath-Base to mathematical instruction tuning using datasets geared toward chain-of-thought \citep{cot}, program-of-thought \citep{pot,pal}, and tool-integrated reasoning \citep{tora} paradigms. The resulting model, DeepSeekMath-Instruct 7B, surpasses all other available 7B models and demonstrates performance on par with leading 70B open-source models with instruction tuning.

In addition, we propose Group Relative Policy Optimization (GRPO), a novel reinforcement learning (RL) approach that modifies the traditional Proximal Policy Optimization (PPO) algorithm \citep{schulman2017proximal}. GRPO eliminates the need for a critic network, instead deriving the baseline directly from grouped sample scores, which results in a significant reduction in computational costs. Leveraging only a selected set of English instruction tuning instances, GRPO delivers significant improvements over the robust DeepSeekMath-Instruct model, elevating performance both in core tasks (GSM8K: 82.9\% $\rightarrow$ 88.2\%, MATH: 46.8\% $\rightarrow$ 51.7\%) and in external mathematical benchmarks (such as CMATH: 84.6\% $\rightarrow$ 88.8\%) during the reinforcement learning stage. We further introduce a cohesive framework that contextualizes various methods, including Rejection Sampling Fine-Tuning (RFT) \citep{yuan2023scaling}, Direct Preference Optimization (DPO) \citep{dpo}, PPO, and GRPO, under the lens of direct or simplified RL strategies. Through this unified perspective, we systematically explore distinctions via comprehensive experiments—such as comparing online and offline training, outcome versus process supervision, and single-turn versus iterative RL—to probe the fundamental drivers underlying these methodologies. Finally, we provide insight into the mechanisms through which RL enhances instruction-tuned models, and offer prospective pathways for realizing more efficient RL approaches within this unifying conceptual scheme.

\subsection{Contributions}
This work advances scalable mathematical pre-training and presents an in-depth examination of reinforcement learning methodologies.

\textbf{Math Pre-Training at Scale}

\begin{itemize}[topsep=0pt]
    \item We demonstrate that the widely available Common Crawl dataset holds significant potential for mathematical research. Utilizing a rigorously crafted data selection framework, we assemble the DeepSeekMath Corpus—a comprehensive dataset containing 120B tokens, curated from mathematically relevant web pages. This resource substantially surpasses previous collections, being nearly sevenfold the size of the Minerva math web dataset \citep{minerva} and approximately nine times larger than OpenWebMath \citep{openwebmath}.

    \item Our foundational model, DeepSeekMath-Base 7B, matches the performance level of Minerva 540B \citep{minerva}, highlighting that scale in model parameters is not the sole determinant of mathematical reasoning performance. Models of more modest scale, when pre-trained on high-caliber datasets, can also yield robust outcomes.

    \item We report insights gained from experiments in mathematical model training, revealing that pre-training on code tasks prior to mathematical content significantly strengthens the model’s proficiency in solving mathematical problems, regardless of whether auxiliary tools are employed. This finding offers a partial resolution to the debated issue: \textit{does code training enhance reasoning abilities?} Our results suggest that, at least in mathematical contexts, the answer is affirmative.

    \item While leveraging arXiv papers for training has become widespread in math-oriented research, our study finds that such data does not lead to discernible gains across the set of mathematical benchmarks utilized in this work.
\end{itemize}

\textbf{Exploration and Analysis of Reinforcement Learning}

\begin{itemize}[topsep=0pt]
    \item We present Group Relative Policy Optimization (GRPO), a novel and resource-efficient reinforcement learning algorithm. In contrast to traditional approaches, GRPO eliminates the need for a critic by deriving the baseline directly from group scores. This innovation leads to a substantial reduction in computational overhead when compared to Proximal Policy Optimization (PPO).
    \item Our experiments reveal that GRPO yields notable improvements in the performance of the instruction-tuned model DeepSeekMath-Instruct using only instruction-tuning data. In addition, we detect marked gains in out-of-domain generalization during reinforcement learning.
    \item We propose a consolidated framework that facilitates the interpretation of methods such as RFT, DPO, PPO, and GRPO. To thoroughly dissect this framework, we perform comprehensive ablations—including comparisons between online and offline learning, outcome versus process-based supervision, as well as single-step and iterative reinforcement learning—to elucidate the foundational components.
    \item Drawing on this unified perspective, we investigate the underlying drivers behind reinforcement learning’s success, and articulate several promising avenues for future work to further enhance reinforcement learning for LLMs.
\end{itemize}

\subsection{Summary of Evaluations and Metrics}

\begin{itemize}[topsep=0pt]
    \item \textbf{English and Chinese Mathematical Reasoning}: We execute a detailed evaluation of our models using a suite of benchmarks for English and Chinese mathematical reasoning, which span tasks from elementary to university level. The English evaluation suite encompasses datasets such as GSM8K \citep{gsm8k}, MATH \citep{MATH}, SAT \citep{llemma}, OCW Courses \citep{minerva}, and MMLU-STEM \citep{mmlu}. The Chinese suite covers MGSM-zh \citep{mgsm}, CMATH \citep{wei2023cmath}, Gaokao-MathCloze \citep{agieval}, and Gaokao-MathQA \citep{agieval}. Assessments are conducted both for the ability to produce fully self-contained written solutions without auxiliary tools and for solving using Python code.

    Within the English benchmarks, DeepSeekMath-Base matches or surpasses the performance of the proprietary Minerva 540B \citep{minerva}, outstripping all known open-source base models (including Mistral 7B \citep{mistral} and Llemma-34B \citep{llemma}), with substantial leads regardless of their pre-training in mathematics. For Chinese datasets, DeepSeekMath-Base exhibits superior results, which may be attributed to our deviation from previous practices \citep{minerva,llemma}—we expand pre-training to include curated, high-quality non-English mathematical data rather than restricting to English only. Following mathematical instruction tuning and reinforcement learning, both DeepSeekMath-Instruct and DeepSeekMath-RL achieve impressive results, notably becoming the first open-source models to attain above 50\% accuracy on the challenging MATH competition dataset.
\end{itemize}

\item \textbf{Formal Mathematics}: We assess DeepSeekMath-Base on the informal-to-formal theorem proving benchmark introduced in \citep{dsp_proof} using miniF2F \citep{minif2f} and Isabelle \citep{isabelle} as the proof assistant. DeepSeekMath-Base achieves notable performance in few-shot autoformalization tasks.
\item \textbf{Natural Language Understanding, Reasoning, and Code}: To thoroughly gauge the model’s capabilities in language understanding, reasoning, and programming, we benchmark DeepSeekMath-Base on the Massive Multitask Language Understanding (MMLU) suite \citep{mmlu}, which spans 57 multiple-choice problems across varied domains, BIG-Bench Hard (BBH) \citep{bbh}, consisting of 23 sophisticated tasks typically requiring multi-step logical reasoning, along with HumanEval \citep{codex} and MBPP \citep{mbpp}, both prominent benchmarks for evaluating code generation abilities. Pre-training on mathematical data proves advantageous for both general language comprehension and logical reasoning.

\section{Math Pre-Training}

\subsection{Data Collection and Decontamination}
This section details our methodology for assembling the DeepSeekMath Corpus from Common Crawl data. As illustrated in Figure \ref{fig:data_collect}, the pipeline operates iteratively to curate a large-scale mathematical text collection, commencing from a carefully selected seed dataset (e.g., a compact, high-quality sample of math-focused materials). Importantly, the method can be generalized to other specialized fields, such as programming.

Our initial step is to select OpenWebMath \citep{openwebmath}, a well-vetted set of mathematical web content, as the foundation for the seed corpus. Utilizing this dataset, we train a fastText model \citep{joulin2016fasttext} designed to identify additional web pages with characteristics similar to those in OpenWebMath. Specifically, we randomly draw 500,000 examples from the seed as positive cases and pair them with 500,000 negative samples from Common Crawl. The training is performed with an open-source implementation\footnote{\url{https://fasttext.cc}} using the following parameters: vector size of 256, learning rate of 0.1, a word n-gram maximum length of 3, a minimum frequency threshold of 3, and 3 epochs. To manage the size of the Common Crawl dataset, we apply URL-based and near-duplicate filtering, narrowing the collection to 40B HTML documents. The trained fastText model is then employed to retrieve mathematical texts from this deduplicated subset. To enhance data quality, candidate documents are ranked by their fastText scores, retaining only the top-scoring entries. The final dataset size is determined through pre-training trials with subsets containing the leading 40B, 80B, 120B, and 160B tokens, with the first cycle focusing on the top 40B tokens.

Following the initial round of data collection, a significant portion of mathematical web content remained unretrieved, primarily due to the limited diversity in the positive samples used to train the fastText model. To remedy this, we incorporate additional mathematical web resources to expand the initial dataset and enhance the fastText model's performance. Our strategy begins with segmenting the complete Common Crawl into non-overlapping domains, where a domain comprises web pages that share the same root URL. We then assess, for each domain, what fraction of web pages were captured during the first pass. Domains from which more than 10\% of web pages were collected are labeled as math-oriented (for example, \url{mathoverflow.net}). Within these domains, we conduct manual annotation of URLs linked to mathematical content (such as \url{mathoverflow.net/questions}). Web pages corresponding to these mathematical URLs but not yet collected are subsequently incorporated into the positive training set. This process allows us to augment the fastText model with a more diverse set of positive examples, improving its recall for mathematical content in later iterations. Ultimately, after four such cycles, the process yields 35.5M mathematical web pages, comprising 120B tokens. Notably, by the fourth cycle, almost 98\% of the pages had already been gathered in the third round, prompting us to conclude data collection.

To ensure no overlap with evaluation benchmarks, we adopt the contamination-filtering procedures of \cite{deepseek-coder}, systematically excluding web pages containing question or answer content from major English mathematical benchmarks—including GSM8K \citep{gsm8k} and MATH \citep{MATH}—as well as from Chinese resources such as CMATH \citep{wei2023cmath} and AGIEval \citep{agieval}. Specifically, if any sequence of 10 consecutive tokens in a web page exactly matches a substring from any benchmark dataset, that page is excluded from the math training corpus. For benchmark excerpts shorter than 10 tokens but at least 3 tokens in length, we similarly apply exact substring filtering to identify and remove contaminated content.

\subsection{Validating the Quality of the DeepSeekMath~Corpus}

We conduct pre-training trials to assess how the DeepSeekMath~Corpus measures up to recently introduced mathematical corpora:

\begin{itemize}[topsep=0pt]
    \item \textbf{MathPile} \citep{mathpile}: a collection drawn from various sources (totaling 8.9B tokens), such as textbooks, Wikipedia, ProofWiki, CommonCrawl, StackExchange, and arXiv—with over 85\% of the data sourced from arXiv;
    \item \textbf{OpenWebMath} \citep{openwebmath}: mathematical data filtered from CommonCrawl, comprising 13.6B tokens in total;
    \item \textbf{Proof-Pile-2} \citep{llemma}: an integrated mathematical corpus, including OpenWebMath, AlgebraicStack (10.3B tokens of math-related code), and arXiv texts (28.0B tokens). For experiments on Proof-Pile-2, we follow the ratio for arXiv:Web:Code as specified by \cite{llemma}, which is 2:4:1.
\end{itemize}

\subsubsection{Training Setting}
\label{sec:quality-policy}
Our approach involves further math-specific training of a 1.3B-parameter, general-purpose language model that utilizes the architecture of the DeepSeek LLMs \citep{deepseek-llm}, referred to herein as DeepSeek-LLM 1.3B. Separate models are fine-tuned on each individual mathematical corpus, with each training process comprising 150B tokens. All model training is conducted with the resource-efficient HAI-LLM framework \citep{haillm}. Consistent with DeepSeek LLMs’ optimization protocols, we employ the AdamW optimizer \citep{adamW}, set to $\beta_1=0.9$, $\beta_2=0.95$, and $\mathrm{weight\_decay}=0.1$. The learning rate follows a multi-step schedule, ramping up over 2,000 warmup steps to its maximum, decreasing to 31.6\% of this peak by the 80\% training mark, and then further dropping to 10.0\% at 90\% of total steps. The maximum learning rate is fixed at 5.3e-4, using batch sizes of 4M tokens and context windows of 4K.

\subsubsection{Evaluation Results}

\textbf{The DeepSeekMath~Corpus stands out for its superior quality, multilingual breadth in mathematical data, and exceptional scale.}

\begin{itemize}[topsep=0pt]
    \item \textbf{High-quality}:
    Downstream evaluation using eight mathematical benchmarks and few-shot chain-of-thought prompting \cite{cot} reveals clear advantages for the model trained on DeepSeekMath~Corpus. As illustrated in Table \ref{tab:corpora_comparison}, this model consistently surpasses others. Additionally, Figure \ref{fig:corpora_comparisons} demonstrates that at the 50B-token mark (representing a full pass over Proof-Pile-2), the model trained with DeepSeekMath Corpus continues to outperform, indicating a higher average dataset quality.
    \item \textbf{Multilingual}:
    With content spanning multiple languages, DeepSeekMath~Corpus predominantly includes English and Chinese materials. Table \ref{tab:corpora_comparison} confirms that training with DeepSeekMath~Corpus improves mathematical reasoning capabilities in both languages. On the other hand, existing corpora, which are overwhelmingly English-focused, yield little improvement—and may even be detrimental—for mathematical reasoning in Chinese.
    \item \textbf{Large-scale}:
    DeepSeekMath~Corpus substantially exceeds the size of prior mathematical datasets. As shown in Figure \ref{fig:corpora_comparisons}, DeepSeek-LLM 1.3B models trained on DeepSeekMath~Corpus display faster and more sustained performance gains. In contrast, the much smaller baseline datasets have already undergone multiple passes, with their associated model improvements stagnating rapidly.
\end{itemize}

\subsection{Training and Evaluating DeepSeekMath-Base 7B}

This section presents DeepSeekMath-Base 7B, a foundational model exhibiting advanced reasoning proficiency, particularly in mathematical contexts. The initial parameters of our model are derived from DeepSeek-Coder-Base-v1.5 7B \citep{deepseek-coder}, after which it undergoes training over 500 billion tokens. The training corpus comprises 56\% DeepSeekMath Corpus, 4\% AlgebraicStack, 10\% arXiv, 20\% Github code, and 10\% natural language material sourced from Common Crawl, incorporating both English and Chinese. Except for adjusting the maximum learning rate to 4.2e-4 and configuring the batch size to 10 million tokens, we adhere to the training procedures outlined in Section \ref{sec:quality-policy}.

A thorough evaluation is performed to analyze the mathematical reasoning and problem-solving capabilities of DeepSeekMath-Base 7B, which includes its proficiency in constructing fully explained solutions autonomously, resolving mathematical questions with auxiliary tools, and conducting rigorous theorem verification. Furthermore, to establish a broader characterization, we also assess the model’s general abilities, such as comprehension of natural language, logical reasoning, and programming competencies.

\paragraph{Mathematical Problem Solving with Step-by-Step Reasoning}

To measure DeepSeekMath-Base's mathematical problem-solving capacity, we utilize few-shot chain-of-thought prompting \citep{cot} across eight distinct English and Chinese benchmarks. These tasks span quantitative reasoning assessments (including GSM8K \citep{gsm8k}, MATH \citep{MATH}, and CMATH \citep{wei2023cmath}) and multiple-choice formats (such as MMLU-STEM \citep{mmlu} and Gaokao-MathQA \citep{agieval}), thereby addressing a wide range of mathematical topics from elementary to collegiate levels.

As detailed in Table \ref{tab:base_math_cot}, DeepSeekMath-Base 7B achieves superior results on all eight evaluated benchmarks when compared with other open-source base models, such as the popular Mistral 7B \citep{mistral} and the newly released Llemma 34B \citep{llemma}, which incorporates math-specific training from Proof-Pile-2 \citep{llemma}. Of particular note, on the MATH benchmark designed for mathematical competition problems, DeepSeekMath-Base outperforms prior open-source base models by more than 10\% in absolute terms, and even surpasses Minerva 540B \citep{minerva}—a much larger closed-source base model based on PaLM \citep{palm} and further specialized on mathematical material—despite being 77 times smaller.

\paragraph{Mathematical Problem Solving with Tool Use} For program-supported mathematical reasoning, we evaluate few-shot program-of-thought prompting \citep{pot,pal} on GSM8K and MATH. Here, models are asked to solve each problem by generating a Python program, making use of libraries like \textit{math} and \textit{sympy} to perform complex calculations. The output produced by running the generated code serves as the model's answer. According to Table \ref{tab:base_math_pot_proof}, DeepSeekMath-Base 7B establishes new state-of-the-art results, outperforming the previous top-performing Llemma 34B.

\paragraph{Formal Mathematics}

Automating formal proof generation is critical for enhancing the rigor, correctness, and productivity of mathematical proof development, a topic that has garnered increasing research attention recently. We assess DeepSeekMath-Base 7B on the informal-to-formal proof task from \citep{dsp_proof}, which requires translating an informal statement, its formalized version, and an informal proof into a formal proof. This is conducted using miniF2F \citep{minif2f}, which is designed for Olympiad-style formal mathematics; for each question, the model produces an Isabelle formal proof with few-shot prompting. Consistent with \cite{dsp_proof}, we first have the model produce a proof sketch and then use the automated prover Sledgehammer \citep{sledgehammer} to complete the proof details. The results in Table \ref{tab:base_math_pot_proof} show that DeepSeekMath-Base 7B delivers robust performance in automating formal proofs.

\paragraph{Natural Language Understanding, Reasoning, and Code}

To gauge natural language understanding, we evaluate on MMLU \citep{mmlu}; for reasoning, we use BBH \citep{bbh}; and for programming abilities, we employ HumanEval \citep{codex} and MBPP \citep{mbpp}. As indicated in Table \ref{tab:understanding_reasoning_code}, DeepSeekMath-Base 7B markedly improves upon its predecessor, DeepSeek-Coder-Base-v1.5 \citep{deepseek-coder}, on MMLU and BBH, demonstrating that focused math training enhances both language understanding and reasoning capacities. Further, the integration of code tokens during ongoing training ensures that DeepSeekMath-Base 7B retains comparable performance to DeepSeek-Coder-Base-v1.5 on both code benchmarks. Across the board, DeepSeekMath-Base 7B substantially exceeds the generalist model Mistral 7B \citep{mistral} in the three reasoning and coding tasks.

\section{Supervised Fine-Tuning}

\subsection{SFT Data Curation}

We assemble an instruction-tuning dataset for mathematics encompassing both English and Chinese language problems drawn from a variety of mathematical disciplines and spanning a range of difficulty. Each problem is coupled with a detailed solution, presented in chain-of-thought (CoT) \citep{cot}, program-of-thought (PoT) \citep{pot,pal}, or tool-integrated reasoning style \citep{tora}. The curated training set comprises 776K examples.

\begin{itemize}[topsep=0pt]
    \item \textbf{English mathematical datasets}: Tool-integrated annotations are provided for items from GSM8K and MATH, while a selected portion of MathInstruct \citep{MathInstruct} and the training partition of Lila-OOD \citep{lila} are included, featuring solutions derived via CoT or PoT approaches. This collection encompasses an extensive range of mathematical areas, including algebra, probability, number theory, calculus, and geometry.
    \item \textbf{Chinese mathematical datasets}: We gather K-12 mathematics questions in Chinese, distributed over 76 distinct sub-domains (such as linear equations), with annotations given in both CoT and tool-integrated reasoning paradigms.
\end{itemize}

\subsection{Training and Evaluating DeepSeekMath-Instruct 7B}

This section details the mathematical instruction-tuning process for DeepSeekMath-Instruct 7B, which is built upon the DeepSeekMath-Base architecture. For training, examples are concatenated in random order until the aggregate reaches a context window of 4K tokens. The training process consists of 500 iterations using a batch size of 256, maintaining a fixed learning rate of 5e-5.

To assess mathematical reasoning capabilities, we evaluate model variants with and without access to external tools across four quantitative reasoning benchmarks in English and Chinese. Our model’s performance is compared with state-of-the-art contemporary systems:

\begin{itemize}[topsep=0pt]
    \item \textbf{Closed-source models}: This comparison set features: (1) members of the GPT family, most notably GPT-4 \citep{gpt4} and GPT-4 Code Interpreter~\footnote{\url{https://openai.com/blog/chatgpt-plugins\#\#code-interpreter}}, (2) Gemini Ultra and Pro \citep{gemini}, (3) Inflection-2 \citep{inflection-2}, (4) Grok-1~\footnote{\url{https://x.ai/model-card}}, as well as recent models introduced by Chinese companies such as (5) Baichuan-3~\footnote{\url{https://www.baichuan-ai.com}}, and (6) the newest GLM-4~\footnote{\url{https://open.bigmodel.cn/dev/api\#glm-4}} from the GLM series \citep{glm}.
\end{itemize}

These models are designed for broad applicability and have typically gone through several alignment stages.

\item \textbf{Open-source models} encompass: general-purpose systems such as (1) DeepSeek-LLM-Chat 67B \citep{deepseek-llm}, (2) Qwen 72B \citep{qwen}, (3) SeaLLM-v2 7B \citep{seallm}, and (4) ChatGLM3 6B \citep{chatglm3}. Additionally, there are models with explicit improvements for mathematical tasks: (5) InternLM2-Math 20B~\footnote{\url{https://github.com/InternLM/InternLM-Math}}, developed from InternLM2 with further math-centric pretraining and subsequent instruction tuning; (6) Math-Shepherd-Mistral 7B, which leverages PPO \citep{schulman2017proximal} training for Mistral 7B \citep{mistral} guided by a reward model focused on process supervision; (7) the WizardMath series \citep{wizardmath}, which enhances mathematical reasoning capabilities in Mistral 7B and Llama-2 70B \citep{llama2} through evolve-instruct—a variant of instruction tuning with machine-generated instructions—and PPO, primarily training on data from GSM8K and MATH; (8) MetaMath 70B \citep{metamath}, where Llama-2 70B is fine-tuned using augmented GSM8K and MATH datasets; (9) ToRA 34B \cite{tora}, resulting from CodeLlama 34B refined for tool-based mathematical reasoning; and (10) MAmmoTH 70B \citep{MathInstruct}, an instruction-tuned variant of Llama-2 70B specialized on MathInstruct.

As detailed in Table \ref{tab:sft_rl_math}, when tool-assisted solutions are prohibited, DeepSeekMath-Instruct 7B excels in stepwise reasoning capabilities. Remarkably, on the advanced MATH benchmark, this model outperforms every open-source alternative and most closed models (such as Inflection-2 and Gemini Pro), surpassing them by a margin of at least 9 percentage points. This advantage holds even in comparison to much larger models (such as Qwen 72B) and those explicitly fine-tuned for mathematical reasoning via reinforcement learning (e.g., WizardMath-v1.1 7B). Although DeepSeekMath-Instruct is on par with proprietary Chinese models like GLM-4 and Baichuan-3 on MATH, it does not reach the results of GPT-4 or Gemini Ultra.

In an evaluation context where both natural language reasoning and the use of external tools are permitted for solving problems, DeepSeekMath-Instruct 7B attains close to 60\% accuracy on MATH, outperforming all other publicly available models. On other standard datasets, its performance rivals that of DeepSeek-LLM-Chat 67B, the former leading model in this category and an order of magnitude larger in size.

\section{Reinforcement Learning}

\subsection{Group Relative Policy Optimization} Reinforcement learning (RL) has demonstrated notable effectiveness in enhancing the mathematical reasoning skills of LLMs following the Supervised Fine-Tuning (SFT) phase \citep{wang2023math,wizardmath}. In this part, we present Group Relative Policy Optimization (GRPO), a novel and resource-efficient RL algorithm.

\subsubsection{From PPO to GRPO} Proximal Policy Optimization (PPO) \citep{schulman2017proximal} is a widely adopted actor-critic reinforcement learning method frequently employed during the RL fine-tuning of LLMs \citep{ouyang2022training}. Its objective is to refine LLMs by optimizing the surrogate function:

\begin{equation}
\footnotesize
    \mathcal{J}_{PPO}(\theta) = \mathbb{E}{[q \sim P(Q), o \sim \pi_{\theta_{old}}(O|q)]} \frac{1}{|o|} \sum_{t=1}^{|o|} \min \left[ \frac{\pi_\theta(o_{t} | q, o_{<t})}{\pi_{\theta_{old}}(o_{t} | q, o_{<t})} A_{t}, \text{clip} \left( \frac{\pi_\theta(o_{t} | q, o_{<t})}{\pi_{\theta_{old}}(o_{t} | q, o_{<t})}, 1 - \varepsilon, 1 + \varepsilon

Here, $\pi_{\theta}$ denotes the updated policy model, while $\pi_{\theta_{old}}$ refers to the previous version of the policy; both $q$ and $o$ stand for questions and responses sampled respectively from the question dataset and from the old policy $\pi_{\theta_{old}}$. The parameter $\varepsilon$ is a clipping hyper-parameter introduced in PPO to promote stable optimization. $A_t$ is the advantage term, calculated using Generalized Advantage Estimation (GAE) \citep{gae}, utilizing the reward sequence $\{r_{\ge t}\}$ and a value function $V_{\psi}$ that is trained during the process. Consequently, PPO requires simultaneous training of a value function alongside the policy. To prevent excessive optimization against the reward model, a standard mitigation involves incorporating a per-token KL divergence penalty—computed against a reference model—into the reward signal at every token \citep{ouyang2022training}, such that

\begin{equation}
    r_{t} = r_\varphi(q, o_{\le t}) - \beta \log\frac{\pi_{\theta}(o_{t}|q, o_{<t})}{\pi_{ref}(o_{t}|q, o_{<t})},
\label{eq:PPO-reward}
\end{equation}

with $r_\varphi$ representing the reward model, $\pi_{ref}$ as the reference (often the initial SFT model), and $\beta$ being the KL penalty coefficient.

Because the value function in PPO typically matches the policy model in complexity and size, its inclusion significantly increases computational and memory requirements. Furthermore, the value function acts as a baseline in advantage calculation for variance reduction during RL, but with LLMs, reward models frequently assign non-zero reward only to the last token of the output. This sparsity presents challenges for training accurate per-token value functions. To overcome these obstacles, we introduce Group Relative Policy Optimization (GRPO), depicted in Figure \ref{fig:grpo}, which removes the need for a separate value function. GRPO instead utilizes the average reward of a set of responses sampled for each question as a baseline. Concretely, for a given question $q$, GRPO generates a group $\{o_1, o_2, \cdots, o_G\}$ of outputs from $\pi_{\theta_{old}}$, and optimizes the new policy with the objective

\begin{equation}
\footnotesize
\begin{split}
    \mathcal{J}_{GRPO}(\theta) &= \mathbb{E}{[q \sim P(Q), \{o_i\}_{i=1}^G \sim \pi_{\theta_{old}}(O|q)]}  \\
    & \frac{1}{G}\sum_{i=1}^G\frac{1}{|o_i|} \sum_{t=1}^{|o_i|} \left\{ \min \left[ \frac{\pi_\theta(o_{i,t} | q, o_{i,<t})}{\pi_{\theta_{old}}(o_{i,t} | q, o_{i,<t})} \hat{A}_{i,t}, \text{clip} \left( \frac{\pi_\theta(o_{i,t} | q, o_{i,<t})}{\pi_{\theta_{old}}(o_{i,t} | q, o_{i,<t})}, 1 - \varepsilon, 1 + \varepsilon \right)  \hat{A}_{i,t} \right] - \beta \mathbb{D}_{KL}\left[\pi_{\theta} || \pi_{ref}\right]\right\} ,
\end{split}
\label{eq:GRPO-obj}
\end{equation}

where $\varepsilon$ and $\beta$ serve as hyper-parameters, and $\hat{A}_{i,t}$ represents the advantage estimated based solely on the relative rewards within each sampled group—this calculation will be elaborated on in the next subsections. This group-relative approach aligns naturally with the way reward models are designed and trained, i.e., through direct comparison of outputs corresponding to the same question. Notably, instead of embedding the KL penalty in the reward, GRPO regularizes the policy update by directly appending the KL divergence between the current and reference policies as a loss term, thereby simplifying the advantage calculation for $\hat{A}_{i,t}$. Unlike the per-token KL penalty seen in (\ref{eq:PPO-reward}), we compute the KL divergence in GRPO using an unbiased estimator proposed in \citep{kl_approx}:

\begin{equation}
\small
    \mathbb{D}_{KL}\

\begin{algorithm}[t]
  \small
  \caption{Iterative Group Relative Policy Optimization}
  \textbf{Input} initial policy model $\pi_{\theta_{\text{init}}}$; reward models $r_\varphi$; task prompts $\mathcal{D}$; 
  hyperparameters $\varepsilon$, $\beta$, $\mu$
  \begin{algorithmic}[1]
    \State Initialize policy model $\pi_\theta \leftarrow \pi_{\theta_{\text{init}}}$
    \For{iteration = 1, \dots, I}
       \State Assign reference model $\pi_{ref} \leftarrow \pi_{\theta}$
      \For{step = 1, \dots, M}
      \State Draw a minibatch $\mathcal{D}_b$ from $\mathcal{D}$
      \State Set old policy model $\pi_{\theta_{old}} \leftarrow \pi_{\theta}$ 
      \State For each prompt $q \in \mathcal{D}_b$, sample $G$ outputs $\{o_i\}_{i=1}^G \sim \pi_{\theta_{old}} (\cdot \mid q)$
      \State For all generated outputs $o_i$, compute associated rewards $\{r_i\}_{i=1}^{G}$ using $r_{\varphi}$
      \State Derive $\hat{A}_{i,t}$ for the $t$-th token in $o_i$ by group-based relative advantage calculation
      \For{GRPO iteration = 1, \dots, $\mu$}
        \State Update policy $\pi_{\theta}$ to maximize the GRPO objective (see Equation \ref{eq:GC-GRPO})
      \EndFor
    \EndFor 
    \State Perform continual training for $r_\varphi$ via replay-based updates
    \EndFor 
  \end{algorithmic}
  \textbf{Output} $\pi_\theta$
  \label{alg:iter-grpo}
\end{algorithm}

\subsubsection{Outcome Supervision RL with GRPO} To define this formally, for any given prompt $q$, we sample a set of $G$ outputs $\{o_1, o_2, \cdots, o_G\}$ from the previous policy version $\pi_{\theta_{old}}$. Each output is evaluated by the reward model, resulting in reward values $\mathbf{r}=\{r_1, r_2, \cdots, r_G\}$. These raw rewards are standardized by subtracting the batch mean and dividing by the batch standard deviation. Under outcome supervision, the normalized reward is delivered only at the terminal state of each output $o_i$; this value is used for all token-level advantages $\hat{A}_{i, t}$, i.e., $\hat{A}_{i, t} = \widetilde{r}_i = \frac{r_i- {\rm mean}(\mathbf{r})}{{\rm std}(\mathbf{r})}$. The policy parameters are then optimized by maximizing the target in equation (\ref{eq:GRPO-obj}).

\subsubsection{Process Supervision RL with GRPO} In contrast, outcome-based supervision supplies feedback solely at sequence completion, which may prove inadequate for intricate mathematical reasoning tasks. Inspired by \cite{wang2023math}, we incorporate process supervision, which allocates rewards at the conclusion of each intermediate reasoning step. Specifically, given prompt $q$ and group of sampled responses $\{o_1, o_2, \cdots, o_G\}$, a process-level reward model is employed to assign reward at every reasoning step, generating $\mathbf{R} = \{ \{r_1^{index(1)},\cdots,r_1^{index(K_1)}\}, \cdots,  \{r_G^{index(1)},\cdots,r_G^{index(K_G)}\} \}$, where $index(j)$ locates the final token of step $j$ and $K_i$ indicates the number of steps in output $i$. These per-step rewards are likewise normalized: $\widetilde{r}_i^{index(j)} = \frac{r_i^{index(j)} - {\rm mean}(\mathbf{R})}{{\rm std}(\mathbf{R})}$. The resulting token-wise advantage $\hat{A}_{i, t}$ aggregates the normalized rewards for subsequent steps, i.e., $\hat{A}_{i, t} = \sum_{index(j) \ge t} \widetilde{r}_i^{index(j)}$. Finally, the model is improved by maximizing the same objective as equation (\ref{eq

\subsubsection{Iterative RL with GRPO} As reinforcement learning progresses, the previously trained reward model may become inadequate for effectively supervising the updated policy model. Thus, we further examine an iterative reinforcement learning framework utilizing GRPO. Illustrated in Algorithm \ref{alg:iter-grpo}, this approach repeatedly regenerates new datasets for the reward model by sampling from the current policy model. The reward model is updated continually, employing a replay buffer that includes 10\% of previous training examples. Subsequently, the policy model serves as the new reference model and is further optimized with the updated reward model.

\subsection{Training and Evaluating DeepSeekMath-RL}

We perform reinforcement learning using DeepSeekMath-Instruct 7B as the policy backbone. The RL training dataset comprises chain-of-thought style questions sourced from GSM8K and MATH within the SFT dataset, amounting to approximately 144K items. Other SFT questions are excluded to specifically assess RL effects on benchmarks without any training data seen during the RL stage. The construction of the reward model’s training set aligns with the methodology from \citep{wang2023math}. Our initial reward model is fine-tuned starting from DeepSeekMath-Base 7B with a learning rate of 2e-5. For the GRPO updates, the policy model uses a learning rate of 1e-6. We use a KL penalty coefficient of 0.04. Each sample consists of $64$ generated outputs per question, with a maximum sequence length set to 1024 and a batch size of 1024 for training. The policy model undergoes a single parameter update after each exploration iteration. DeepSeekMath-RL 7B is evaluated on the same benchmarks as DeepSeekMath-Instruct 7B. Specifically, GSM8K and MATH datasets that require chain-of-thought solutions are categorized as in-domain tasks for DeepSeekMath-RL 7B, while the remaining benchmarks are treated as out-of-domain evaluation tasks.

Table \ref{tab:sft_rl_math} presents comparative results for both open- and closed-source models, incorporating chain-of-thought and tool-assisted reasoning on English and Chinese datasets. Our analysis reveals the following: (1) DeepSeekMath-RL 7B achieves 88.2\% accuracy on GSM8K and 51.7\% on MATH under chain-of-thought settings. These outcomes outperform all open-source competitors in the 7B–70B parameter range, as well as surpass most closed-source alternatives. (2) Notably, DeepSeekMath-RL 7B’s training is limited exclusively to chain-of-thought instruction-tuning data from GSM8K and MATH, initialized from DeepSeekMath-Instruct 7B. Despite the narrowly focused training corpus, it demonstrates consistent improvements over DeepSeekMath-Instruct 7B across every measured metric, underscoring the value added by reinforcement learning.

\section{Discussion}
In this section, we detail key insights gleaned from our pre-training and reinforcement learning studies.

\subsection{Lessons Learnt in Pre-Training}

We begin by summarizing observations from our pre-training efforts. Unless indicated otherwise, we reference the training regime described in Section \ref{sec:quality-policy}. Additionally, throughout this section, the term DeepSeekMath Corpus refers to an 89B-token collection derived from the second cycle of data aggregation.

\subsubsection{Code Training Benefits Mathematical Reasoning}

The hypothesis that code-related pre-training boosts reasoning capabilities has garnered much attention, yet remains largely unsubstantiated. Our investigation provides evidence within the mathematical sphere: exposure to code during training enhances models’ mathematical reasoning skills, independent of whether external tools are leveraged.

To evaluate the influence of code pre-training on mathematical reasoning, we designed two experimental regimes: a two-stage and a single-stage training approach.

\textbf{Two-Stage Training}

\begin{itemize}[topsep=0pt]
    \item \textbf{Code Training for 400B Tokens $\rightarrow$ Math Training for 150B Tokens}: The DeepSeek-LLM 1.3B model is initially trained on 400B tokens derived from code, subsequently followed by an additional 150B tokens focusing on mathematical content.
    \item \textbf{General Training for 400B Tokens $\rightarrow$ Math Training for 150B Tokens}: As a baseline, we also carry out training using general corpus tokens (sampled from a comprehensive dataset compiled by DeepSeek-AI) rather than code tokens in the initial stage, to directly evaluate the comparative influence of code tokens and general language tokens on the development of mathematical reasoning abilities.
\end{itemize}

\textbf{One-Stage Training}

\begin{itemize}[topsep=0pt]
    \item \textbf{Math Training for 150B Tokens}: Here, DeepSeek-LLM 1.3B is trained exclusively on 150B math tokens in a single stage;
    \item \textbf{Training on a mixture of 400B Code Tokens and 150B Math Tokens}: Given that math training following a dedicated code pretraining stage negatively impacts code performance, we assess if simultaneously training on a combined set of 400B code tokens and 150B math tokens might reinforce mathematical reasoning while alleviating catastrophic forgetting of coding proficiency.
\end{itemize}

\paragraph{Results}
Performance outcomes for these various training regimens are summarized in Table \ref{tab:code-math} and Table \ref{tab:code-math-general-eval}.

Introducing code tokens during training has a positive impact on program-augmented mathematical problem-solving, across both sequential (two-stage) and simultaneous (one-stage) training approaches. Specifically, Table \ref{tab:code-math} shows that, within a two-stage schedule, a code-pretraining phase meaningfully boosts performance on both GSM8K and MATH datasets when leveraging Python for solution synthesis. Incorporating a subsequent stage of mathematical training drives additional advances. Notably, in the one-stage setting, integrating code and mathematical tokens during joint training substantially counters the catastrophic forgetting typically observed after sequential training. This approach also delivers combined benefits for coding skills (Table \ref{tab:code-math-general-eval}) and program-supported mathematical tasks (Table \ref{tab:code-math}).

The inclusion of code tokens enhances mathematical reasoning even without relying on external computational tools. In two-stage training, pretraining on code results in measurable improvements and further accelerates progress in the later math-focused phase, ultimately yielding the highest overall performance. However, concurrently training on both code and math tokens in a single stage leads to diminished results in mathematical reasoning tasks that do not involve tool use. This reduction could stem from DeepSeek-LLM 1.3B’s constrained capacity, which may be insufficient to thoroughly internalize both code and mathematical information during concurrent exposure.

\subsubsection{ArXiv Papers Appear Unhelpful for Enhancing Mathematical Reasoning}

ArXiv papers have become a standard part of math pre-training datasets \citep{minerva,gpt-f,llemma,mathpile}, yet few studies have systematically investigated their actual influence on mathematical reasoning capabilities. Somewhat surprisingly, our empirical findings indicate that arXiv content does not significantly enhance mathematical reasoning performance. We evaluated this effect across models of varying sizes, specifically DeepSeek-LLM 1.3B and DeepSeek-Coder-Base-v1.5 7B \citep{deepseek-coder}, using distinct arXiv datasets prepared through different data processing strategies:

\begin{itemize}[topsep=0pt]
    \item \textbf{MathPile} \citep{mathpile}: an 8.9B-token dataset constructed using extensive cleaning and heuristic-based filtering, in which more than 85\% of tokens originate from scientific arXiv documents;
    \item \textbf{ArXiv-RedPajama} \citep{redpajama}: a collection of the entire arXiv in LaTeX format, stripped of preambles, comments, user-defined macros, and bibliography sections, comprising 28.0B tokens.
\end{itemize}

Our methodology involved separately pre-training DeepSeek-LLM 1.3B on 150B tokens and DeepSeek-Coder-Base-v1.5 7B on 40B tokens drawn exclusively from each arXiv dataset. Consistently, the results suggest that arXiv-derived data does not yield improvements in mathematical reasoning proficiency. Training models solely on arXiv materials did not result in marked performance gains, and occasionally led to declines across several math evaluation benchmarks with varied levels of challenge, such as quantitative reasoning tasks (GSM8K, MATH; Table \ref{tab:arxiv-cot}), multiple-choice STEM benchmarks (MMLU-STEM; Table \ref{tab:arxiv-cot}), and tasks in formal mathematics (miniF2F; Table \ref{tab:arxiv_atp}).

Nonetheless, this assessment comes with important caveats and should not be viewed as definitive. Specifically, we did not address:

\begin{itemize}[topsep=0pt]
    \item The effects of arXiv-sourced data on other mathematical applications absent from this work, such as translating formal mathematical statements or proofs into informal expositions;
    \item The outcomes of combining arXiv material with additional data types;
    \item The potential for observable benefits from arXiv papers as model scale further increases.
\end{itemize}

Therefore, more comprehensive studies are warranted, which we intend to pursue in future research.

\subsection{Perspectives on Reinforcement Learning}
\subsubsection{A Move Towards a Unified Training Framework} Here, we introduce a unified conceptual framework to interpret diverse learning approaches, including SFT, RFT, DPO, PPO, and GRPO, and experimentally probe key aspects of this unified perspective. In most cases, the gradient with respect to model parameters $\theta$ for these learning schemes may be represented as follows:

\begin{equation}
    \nabla_{\theta}\mathcal{J}_{\textcolor{red}{\mathcal{A}}}(\theta) = \mathbb{E}[\underbrace{(q,o) \sim \textcolor{red}{\mathcal{D}}}_{Data \ Source}]\left( \frac{1}{|o|} \sum_{t=1}^{|o|}  \underbrace{GC_{{\mathcal{A}}}(q, o, t, \textcolor{red}{\pi_{{rf}}})}_{Gradient \ Coefficient}  \nabla_{\theta}\log \pi_{\theta}(o_t | q, o_{<t})\right).
\label{eq:objective}
\end{equation}

The framework is built around three central elements: (1) the \textit{Data Source $\mathcal{D}$}, responsible for supplying training examples; (2) the \textit{Reward Function $\pi_{{rf}}$}, which provides evaluative feedback or reward signals during learning; and (3) the \textit{Algorithm $\mathcal{A}$}, tasked with integrating the data and reward information to produce the gradient coefficient $GC$, which governs how each example influences parameter updates through reinforcement or penalization. Employing this unified formalism, we consider the following noteworthy algorithms:

\begin{itemize}
    \item  \textbf{Supervised Fine-tuning (SFT)}: SFT involves adjusting a pretrained model using labeled data curated by humans.
    \item \textbf{Rejection Sampling Fine-tuning (RFT)}: RFT enhances the SFT-trained model with additional fine-tuning on model outputs—generated for SFT queries—which have been filtered according to answer validity.
    \item \textbf{Direct Preference Optimization (DPO)}: DPO further adapts the SFT model by leveraging preference data and employing a pairwise DPO loss on extended sets of generated outputs.
    \item \textbf{Online Rejection Sampling Fine-tuning (Online RFT)}: Unlike traditional RFT, Online RFT starts from the SFT policy and iteratively fine-tunes using new outputs continuously sampled from the current policy.
    \item \textbf{PPO/GRPO}: In these methods, the SFT-trained model is used to initialize the policy, which is then updated and improved with real-time samples via reinforcement signals.
\end{itemize}

Table \ref{tab:data_policy} provides a comparative summary of the methodological components discussed. For an in-depth explanation of the derivation procedures, please consult Appendix \ref{app:analysis-rl}.

\paragraph{Observation about Data Source} We categorize the sources of training data into online and offline sampling. In the online sampling approach, the dataset is generated by the current training policy model as it explores the environment in real time. Conversely, offline sampling draws from outputs produced by the original SFT model prior to policy updates. RFT and DPO exemplify the offline data usage paradigm, whereas Online RFT and GRPO implement an online data gathering strategy.

According to the results presented in Figure \ref{fig:rl-analysis}, Online RFT consistently achieves superior results compared to RFT across both evaluation benchmarks. In the early training epochs, Online RFT performs on par with RFT; however, it surpasses RFT as training progresses, thereby evidencing the effectiveness of online sampling. This trend aligns with expectations, given that the actor and SFT model share similar behaviors initially, resulting in training samples with minor discrepancies. As training advances, divergence between the actor's outputs and those of the SFT model increases, amplifying the benefits conferred by online, dynamically-generated data.

\paragraph{Observation about Gradient Coefficient} The manner in which reward signals are translated into gradient coefficients fundamentally shapes model parameter updates. Our experiments distinguish between two types of reward functions: ‘Rule’ and ‘Model.’ The ‘Rule’ function assesses the quality of responses by verifying answer correctness, while the ‘Model’ function relies on a trained reward model, itself supervised by rule-based judgements, to assign response scores. Equations \ref{eq:GC-RFT} and \ref{eq:GC-GRPO} elucidate a principal distinction: GRPO uniquely incorporates reward magnitudes into the gradient scaling, allowing the algorithm to proportionally reinforce or penalize responses. By contrast, Online RFT does not implement graded penalization; all correct responses receive uniform positive updates, while incorrect responses are not explicitly penalized.

As depicted in Figure \ref{fig:rl-analysis}, GRPO outperforms online RFT, underlining the effectiveness of modifying the positive and negative gradient coefficients. Moreover, GRPO+PS achieves better results than GRPO+OS, emphasizing the advantage conferred by utilizing more detailed, step-sensitive gradient coefficients. Additionally, we investigate the effect of iterative RL and, for this purpose, implement two rounds of iteration in our experiments. Figure \ref{fig:iter-rl} reveals that iterative RL substantially boosts model performance, with notable gains especially observable during the first iteration.

\subsubsection{Why RL Works?} This work applies reinforcement learning to a subset of instruction tuning data and observes that it yields considerable improvements over models trained solely with instruction tuning. To delve into the mechanisms underlying this success, we compare the Pass@K and Maj@K accuracies for both the Instruct and RL-trained models across two benchmark datasets. Figure \ref{fig:maj-pass} illustrates that RL raises Maj@K scores without corresponding improvements in Pass@K. This suggests that reinforcement learning enhances overall performance by making the output distribution more reliable, specifically by \textbf{elevating the likelihood of the correct answer appearing among the TopK candidates rather than strengthening the model’s core capabilities.} Likewise, \citep{wang2023making} identified a \textbf{misalignment problem} in SFT models when handling reasoning tasks and reported that incorporating preference alignment strategies \citep{yuan2023rrhf,song2023preference,wang2023making} can significantly bolster the reasoning abilities of these models.

\subsubsection{How to Achieve More Effective RL?}
\label{sec:effec_rl}
Our results show that reinforcement learning is highly effective for mathematical reasoning scenarios. We further propose a unified perspective that accommodates a range of widely-used training approaches, interpreting all such methods as direct or approximated forms of RL. As outlined in Equation \ref{eq:objective}, these strategies share three foundational aspects: Data Source, Algorithm, and Reward Function. We discuss promising avenues for future investigation along these dimensions.

\paragraph{Data Source} The data source constitutes the foundation for all training approaches. In the RL framework, we specifically refer to the data source as consisting of unlabeled questions coupled with responses sampled from the policy model. In our study, we employ questions derived from the instruction tuning dataset and sample outputs using a basic nucleus sampling scheme. This methodological choice might partially explain why our RL method leads only to improvements in Maj@K. For future work, we plan to test the RL pipeline on prompts from broader distributions and pair it with \textbf{more advanced sampling (decoding) mechanisms}, such as those informed by tree-search paradigms \citep{yao2023tree}. Additionally, the application of \textbf{efficient inference methodologies} \citep{xia-etal-2023-speculative,leviathan2023fast,kwon2023efficient,xia2024unlocking}, which directly affect how effectively policy models can explore, will also be central to advancing this research.

\paragraph{Algorithms} Algorithms utilize both data and reward signals to adjust the gradient coefficient, which in turn modifies the model parameters. According to Equation \ref{eq:objective}, contemporary approaches predominantly \textbf{TRUST} the feedback from the reward function to adjust the conditional likelihood of a given token. Nonetheless, the reliability of the reward signal cannot always be guaranteed, particularly in highly intricate tasks. For instance, the PRM800K datasets \citep{lightman2023let}, though meticulously labeled by skilled annotators, still possess an estimated 20\% rate of annotation errors\footnote{\url{https://github.com/openai/prm800k/issues/12\#issuecomment-1728491852}}. In light of this, our focus is directed toward reinforcement learning algorithms that exhibit robustness to imperfect or noisy reward signals. We contend that \textbf{WEAK-TO-STRONG} \citep{burns2023weak} alignment strategies may fundamentally transform the landscape of learning algorithms.

\paragraph{Reward Function} The reward function serves as the primary provider of feedback during training. In reinforcement learning, this is commonly implemented as a neural reward model. We identify three crucial avenues for developing reward models: 1) \textbf{Strategies to improve the reward model's generalization.} The reward model needs to generalize effectively to out-of-distribution inputs and sophisticated outputs; if not, reinforcement learning might merely reinforce the status quo in LLM distributions rather than drive true enhancement; 2) \textbf{Techniques to quantify and utilize the uncertainty in reward models.} Properly modeled uncertainty could mediate between underperforming reward models and weak-to-strong learning frameworks; 3) \textbf{Efficient construction of high-fidelity process reward models} capable of delivering nuanced signals to guide reasoning tasks \citep{lightman2023let,wang2023math}.

\section{Conclusion, Limitation, and Future Work}

This paper introduces DeepSeekMath, which surpasses all open-source alternatives on the competition-level MATH benchmark and attains performance levels close to proprietary systems. DeepSeekMath~is based on DeepSeek-Coder-v1.5 7B, having undergone continual pre-training with 500B tokens, of which 120B are mathematics-focused and collected from Common Crawl. Our detailed ablation analyses indicate that web sources provide rich, high-quality mathematical data, whereas arXiv does not yield as much benefit as anticipated. We propose Group Relative Policy Optimization (GRPO)—a modification of Proximal Policy Optimization (PPO)—which enhances mathematical reasoning ability while consuming less memory. Experimental findings confirm that GRPO delivers consistent improvements, even when DeepSeekMath-Instruct 7B achieves strong benchmark results. Furthermore, we establish a unified framework for understanding diverse methodologies and highlight prospective avenues for advancing reinforcement learning techniques.

Despite achieving robust results on quantitative reasoning evaluations, DeepSeekMath~shows relatively limited proficiency in geometry and theorem-proving compared to closed-source models. Specifically, in preliminary assessments, the model struggled with triangle and ellipse problems, suggesting possible biases in pre-training and fine-tuning data selection. Moreover, DeepSeekMath~lags behind GPT-4 in terms of few-shot learning, likely due to its smaller parameter count. While GPT-4 exhibits gains from few-shot prompts, DeepSeekMath~performs similarly under zero-shot and few-shot setups. Future efforts will concentrate on refining our data curation pipeline to assemble a higher-quality pre-training dataset. We also plan to investigate the promising directions presented in Section \ref{sec:effec_rl} to realize more efficient reinforcement learning approaches for LLMs.

\end{document}